\section{Methodology}
\label{sec:methodology}

\note{\textbf{This section is a placement, not new content.} It holds the
measurement paragraph that sat above the abstract in the outline, which is
methodology and belongs here. The text below is that paragraph verbatim.}

Metric is how many paco-specific newly-needed tactics there were before, and how
many of that kind for the new library there are now. We should measure this in
lines and characters of ltac for each of these as well.

We also want total lines of proof and definition, as well as characters of proof
and definition.

The three checkpoints are: paco version, old port, new port.

\note{\textbf{What the released measurement actually does}, for whatever of this
you want to keep. \emph{Checkpoints:} \cpaco{} is ITrees on paco; \cold{} is the
port against the released library; \cnew{} is the port against the rewritten one.
Two scopes are reported separately, \code{theories/} alone and
\code{theories/}~+~\code{extra/}, because they move in opposite directions on
tactic surface. \emph{Tactic counting} uses three conventions: the unit is the
name a user types (\code{step} and \code{step in} are one name); a private helper
folds into the name it serves; a name defined in several files counts once, with
every copy counted in the line and character totals. Membership is curated by
hand, not grepped, because a name-based scan was wrong three times in this work.
\emph{Lines} are counted comment-free with whitespace stripped; blank lines do
not count. \emph{Note that \code{Ltac} is specification for the classifier}, so
tactic definitions are counted inside the definition columns. Several claims in
\Cref{sec:initial} depend on separating them out.}

\note{\textbf{Build time is measured and is not in the outline.} Consider a
subsection here for the protocol, since it changed the answer three times:
\code{make TIMED=1} leaves timing artifacts in one tree and not another; blocked
ordering plus machine warm-up manufactures a difference that is not real, so runs
must interleave; and a single run cannot resolve a two-second effect on a
ninety-second build. Results are in \Cref{tab:timing}.}
